% !TeX root = ../../Book4.tex
% 纲要标题：### F.3 严格n范畴
% 纲要标签：b4:sec:F03
% ---
% 纲要细目（写作范围）：
% <!-- 短节：不设小节 -->

\section{严格n范畴}
\label{b4:sec:F03}

在严格2范畴中，对象之间的 Hom 是范畴，因而其中还有态射。继续这一做法，可以让 Hom 取值于严格2范畴、严格3范畴，依次增加箭头的层次。\cref{b4:def:enriched-category}把这种内部复合写成富集数据；这里所用的张量是笛卡尔积。

\begin{definition}\label{b4:def:F-strict-n-category}
对非负整数 \(n\)，以下同时递归定义小严格 \(n\) 范畴及其严格函子。严格0范畴是集合，严格0函子是集合映射；它们组成范畴 \(0\text{-}\mathbf{Cat}=\mathbf{Set}\)，其积为集合积，终对象为单点集 \(\mathbf1_0\)。

设 \(n\ge1\)，且由小严格 \((n-1)\) 范畴及严格 \((n-1)\) 函子组成的范畴 \((n-1)\text{-}\mathbf{Cat}\) 已有笛卡尔积及终对象 \(\mathbf1_{n-1}\)。给定对象集合 \(\operatorname{Ob}\mathcal C\)，每对对象 \(X,Y\) 对应一个小严格 \((n-1)\) 范畴 \(\operatorname{Hom}_{\mathcal C}(X,Y)\)，以及严格 \((n-1)\) 函子
\[
\begin{gathered}
m_{X,Y,Z}:\operatorname{Hom}_{\mathcal C}(Y,Z)
\times\operatorname{Hom}_{\mathcal C}(X,Y)
\longrightarrow\operatorname{Hom}_{\mathcal C}(X,Z),\\
j_X:\mathbf1_{n-1}\longrightarrow\operatorname{Hom}_{\mathcal C}(X,X).
\end{gathered}
\]
要求下列等式作为严格 \((n-1)\) 函子等式成立，其中积的括号及终对象因子按典范同构识别：
\[
\begin{aligned}
m_{W,X,Z}\circ(m_{X,Y,Z}\times1)
&=m_{W,Y,Z}\circ(1\times m_{W,X,Y}),\\
m_{X,Y,Y}\circ(j_Y\times1)&=1,\\
m_{X,X,Y}\circ(1\times j_X)&=1.
\end{aligned}
\]
这些数据组成一个小\term[yange-n-fanchou]{严格\(n\)范畴}[strict \(n\)-category]，也就是富集于
\[
((n-1)\text{-}\mathbf{Cat},\times,\mathbf1_{n-1})
\]
的范畴。

对两个这样的小严格 \(n\) 范畴 \(\mathcal C,\mathcal D\)，给定对象映射 \(X\mapsto FX\) 及严格 \((n-1)\) 函子
\[
F_{X,Y}:\operatorname{Hom}_{\mathcal C}(X,Y)
\longrightarrow\operatorname{Hom}_{\mathcal D}(FX,FY).
\]
若它们满足
\[
\begin{aligned}
F_{X,Z}\circ m^{\mathcal C}_{X,Y,Z}
&=m^{\mathcal D}_{FX,FY,FZ}\circ(F_{Y,Z}\times F_{X,Y}),\\
F_{X,X}\circ j^{\mathcal C}_X&=j^{\mathcal D}_{FX},
\end{aligned}
\]
则称这些数据为\term[yange-n-hanzi]{严格\(n\)函子}[strict \(n\)-functor]。小严格 \(n\) 范畴及这些函子组成的范畴记为 \(n\text{-}\mathbf{Cat}\)。
\end{definition}
\symidx{\(n\text{-}\mathbf{Cat}\)：小严格 \(n\) 范畴及严格函子的范畴}

严格函子的复合逐层进行：对象映射作通常复合，Hom 上规定 \((GF)_{X,Y}=G_{FX,FY}\circ F_{X,Y}\)。两个函子各自保持复合与单位，把对应的等式接续便得到 \(GF\) 的同样性质；恒等函子在每个 Hom 上取恒等。因此这确实给出范畴 \(n\text{-}\mathbf{Cat}\)。

\ProseParagraphBreak

递归所需的积也逐层构造。严格 \(n\) 范畴 \(\mathcal C\times\mathcal D\) 的对象是对象对，其 Hom 为
\[
\operatorname{Hom}_{\mathcal C\times\mathcal D}\bigl((X,U),(Y,V)\bigr)
=\operatorname{Hom}_{\mathcal C}(X,Y)\times\operatorname{Hom}_{\mathcal D}(U,V).
\]
复合与单位按两个分量给出。投影是严格函子；一对到 \(\mathcal C,\mathcal D\) 的严格函子在对象和各 Hom 上唯一确定到这个积的严格函子，故满足积的泛性质。终对象 \(\mathbf1_n\) 有一个对象，其 Hom 为 \(\mathbf1_{n-1}\)，复合与单位是终对象所决定的唯一函子。每个小严格 \(n\) 范畴到它也恰有一个严格函子。

\ProseParagraphBreak

积的重新分组与复合的严格结合是两件事。通常的积构造把 \((\mathcal A\times\mathcal B)\times\mathcal D\) 的对象写成 \(((a,b),d)\)，把 \(\mathcal A\times(\mathcal B\times\mathcal D)\) 的对象写成 \((a,(b,d))\)；逐层重新分组给出一个固定的典范同构。定义中的识别只是用它把两种复合函子的源对齐；单位律同样先把 \(\mathbf1_{n-1}\times\mathcal A\) 或 \(\mathcal A\times\mathbf1_{n-1}\) 识别为 \(\mathcal A\)。识别以后，两种复合在对象和各层态射上的值都必须相等。例如依次可复合的1态射仍满足 \((h*g)*f=h*(g*f)\)，没有再指定一张连接这两个复合的非平凡高层态射。\secref{b4:sec:F02}中的双范畴则把这样的结合约束另列为数据。

\ProseParagraphBreak

当 \(n=1\) 时，Hom 对象是集合，\(m\) 是集合映射，\(j_X\) 选出恒等态射；公理就是普通小范畴的结合律和单位律，严格1函子就是普通函子。当 \(n=2\) 时，Hom 对象是小范畴，\(m\) 是函子，\(j_X\) 选出单位1态射及它的恒等2态射；于是恢复\cref{b4:def:F-strict-two-category}中的小严格2范畴。严格2函子在对象、1态射和2态射上给出映射，并严格保持两种复合与两层单位。

\ProseParagraphBreak

当 \(n=3\) 时，固定 \(\mathcal C\) 中的对象 \(X,Y\)，令
\[
\mathcal H=\operatorname{Hom}_{\mathcal C}(X,Y).
\]
这是一个严格2范畴，它的对象 \(f,g:X\rightarrow Y\) 是 \(\mathcal C\) 的1态射。再固定 \(f,g\)，则 \(\operatorname{Hom}_{\mathcal H}(f,g)\) 是普通范畴：其中的对象 \(\alpha,\beta:f\Rightarrow g\) 是 \(\mathcal C\) 的2态射，而其中的态射 \(\Theta:\alpha\rightarrow\beta\) 是 \(\mathcal C\) 的3态射。这里把 \(\Theta\) 写成普通箭头，是因为它正在这个普通 Hom 范畴中读取；它在 \(\mathcal C\) 中的层次由两重 Hom 确定。

\ProseParagraphBreak

固定 \(f,g\) 后，3态射可在 \(\operatorname{Hom}_{\mathcal H}(f,g)\) 中复合，记为 \(\Theta'\cdot\Theta\)。若再取 \(\mathcal H\) 的对象 \(h\)，它的复合函子为
\[
\operatorname{Hom}_{\mathcal H}(g,h)\times
\operatorname{Hom}_{\mathcal H}(f,g)
\longrightarrow\operatorname{Hom}_{\mathcal H}(f,h).
\]
记这个函子的值为 \(*_{\mathcal H}\)。它把2态射 \(\alpha:f\Rightarrow g\)、\(\beta:g\Rightarrow h\) 送到2态射 \(\beta*_{\mathcal H}\alpha:f\Rightarrow h\)，也把它们之间的3态射送到相应复合之间的3态射。例如，取 \(\alpha_0,\alpha_1,\alpha_2\in\operatorname{Hom}_{\mathcal H}(f,g)\) 与 \(\beta_0,\beta_1,\beta_2\in\operatorname{Hom}_{\mathcal H}(g,h)\)，以及 \(\Theta_i:\alpha_{i-1}\rightarrow\alpha_i\)、\(\Psi_i:\beta_{i-1}\rightarrow\beta_i\)（\(i=1,2\)），函子性给出
\[
(\Psi_2\cdot\Psi_1)*_{\mathcal H}(\Theta_2\cdot\Theta_1)
=(\Psi_2*_{\mathcal H}\Theta_2)\cdot
 (\Psi_1*_{\mathcal H}\Theta_1).
\]
两边都是从 \(\beta_0*_{\mathcal H}\alpha_0\) 到 \(\beta_2*_{\mathcal H}\alpha_2\) 的3态射。这就是\cref{b4:prop:F-interchange-law}在 \(\mathcal H\) 中的等式；其中的箭头比 \(\mathcal C\) 中相应的箭头低一层。

\ProseParagraphBreak

还可以让 \(\mathcal C\) 中的1态射复合。此时
\[
m_{X,Y,Z}:\operatorname{Hom}_{\mathcal C}(Y,Z)\times
\operatorname{Hom}_{\mathcal C}(X,Y)
\longrightarrow\operatorname{Hom}_{\mathcal C}(X,Z)
\]
是严格2函子。它在对象上把 \((g,f)\) 送到1态射 \(g*f\)；在这两个 Hom 中的1态射上，把2态射对 \((\beta:g\Rightarrow g',\alpha:f\Rightarrow f')\) 送到2态射 \(\beta*\alpha:g*f\Rightarrow g'*f'\)。若还有2态射 \(\alpha':f\Rightarrow f'\)、\(\beta':g\Rightarrow g'\)，则在更高一层，它把3态射对 \((\Psi:\beta\rightarrow\beta',\Theta:\alpha\rightarrow\alpha')\) 送到3态射 \(\Psi*\Theta:\beta*\alpha\rightarrow\beta'*\alpha'\)。严格2函子性要求它保持输入的两个 Hom 严格2范畴中的1态射复合，以及2态射的水平、垂直复合；按 \(\mathcal C\) 的层数计，这些分别是2态射的复合和3态射的两种复合，相应的单位也必须保持。因此先在输入的两个分量中复合，再应用 \(m_{X,Y,Z}\)，与先应用它再复合所得的态射相等。递归定义把这几种复合及其交换律一起规定下来，而不是只要求1态射可以复合。

\ProseParagraphBreak

对一般 \(n\ge1\)，同样逐层读取：最外层 Hom 中的对象是1态射；若 \(n\ge2\)，Hom 内的1态射是2态射，依此直到 \(n\) 态射。若改为弱 \(n\) 范畴，结合与单位等式通常要由指定的可逆高层态射或等价来比较，这些比较又须满足更高层的相容条件；\secref{b4:sec:F02}中双范畴的五边形和三角正是最初的例子。一般弱 \(n\) 范畴需要同时处理这些相干数据，不能仅把上述富集定义中的严格等号逐个换成同构。下一节的拟范畴采用 \((\infty,1)\)-范畴模型，其中二阶及以上的比较在同伦意义下可逆；本节的严格 \(n\) 范畴并未要求这些高层态射可逆。因此拟范畴并不是仅将这里的层数延伸到无穷再放松结合律。
